\section{Additional Considerations, Abbreviations, and Amendment History}

\subsection{Additional Considerations}
\draftinstr{Follow the Phase 2 publication section
\texttt{final-protocol/publication/sections/sec-11-additional.tex} and the Phase 1
model file \texttt{trial-protocol/draft-protocol/sections/sec-11-additional.tex}.
Describe the cross-jurisdictional SaMD and robotically assisted surgery tracks (EU
MDR plus AI requirements, NMPA, PMDA, MHRA, Health Canada, TGA) onto which the
common technical spine maps without redesign, anchored to the consensus standards
\cite{iec8060127,asmevv40,nasastd7009,ieee7009}, and summarize them in the
full-width \textbf{Table~\ref{tab:jurisdictions}}. Describe the federated multicenter
learning model and the H.R. 9510 VVUQ alignment \cite{kawchak2026hr9510}. Source from
\texttt{trial-protocol/nih-protocol/09\_additional\_considerations\_abbreviations\_amendment\_history.md}
and the ten-gate logic from \texttt{trial-protocol/inputs/auto-bill-02/final-bill/main.tex}.}

\draftinstr{Render \texttt{trial-phase-2/mermaid/fig-22-vvuq-ten-gate.md} as a TikZ
\texttt{mermaidfig} (this is \textbf{Figure 22}): verification before generation -
every proposed robot-patient command passes a ten-gate assurance suite (fourteen
external standards plus two clinical baselines) with epistemic and aleatory
uncertainty bounds before the gate surface returns ACCEPT, BLOCK, or ESCALATE.}

\subsection{Protocol Amendment History}
\draftinstr{Build the amendment-history table (\textbf{Table~\ref{tab:amend}}) with
the Phase 1 predicate row (v1.0.0, June 20 2026, the first-in-human single-arm
combined IND/IDE that established the RP2D and device feasibility
\cite{kawchak2026phase1}) and this protocol row (v1.1.0, June 23 2026, the Phase 2
multicenter randomized controlled design with the upgraded Phase 0 gate, single-IRB
governance, the co-investment model and capital firewall, and the confirmatory
group-sequential framework).}

\subsection{Abbreviations}
\draftinstr{Build the two-column-pair abbreviations table
(\textbf{Table~\ref{tab:abbrev}}) resolving every abbreviation used across the
protocol (AE, BICR, CONSORT, CRF, CRO, ctDNA, CTCAE, DSMB, ECOG, EORTC, FDA, GCP,
GJ, HJ, HPB, HR, ICH, IDE, IND, ISGPS, ISM, ITT, KRAS, LLM, MDR, MHRA, mITT, MPR,
MTD, NMPA, OS, PACIF, PDAC, PFS, PJ, PMDA, PP, PV, QoL, R0, RCT, RP2D, SAE, SaMD,
sIRB, SMA, SMV, TGA, USL, VVUQ).}
